\chapter{Experimental Results}

\section{Final Neural Network Performance}

The final neural network results on the held-out test set are reported in Table~\ref{tab:nn-final-metrics}.

\begin{table}[H]
    \centering
    \begin{tabular}{lr}
        \toprule
        Metric             & Value    \\
        \midrule
        MAE                & 0.04290  \\
        RMSE               & 0.06208  \\
        $R^2$              & 0.999993 \\
        MAPE, prices $> 1$ & 0.4803\% \\
        \bottomrule
    \end{tabular}
    \caption{Final neural network metrics against analytical Black-Scholes prices.}
    \label{tab:nn-final-metrics}
\end{table}

The results show that the feed-forward neural network with moneyness as an engineered feature and SiLU hidden activations approximates the Black-Scholes pricing function with very high accuracy on the sampled parameter domain.
The coefficient of determination is close to one, while MAE and RMSE are small in absolute price units.

\begin{figure}[H]
    \centering
    \maybeincludegraphics[width=0.72\linewidth]{../results/final/figures/loss.png}
    \caption{Training loss and validation loss for the final experiment.}
    \label{fig:final-loss}
\end{figure}

Figure~\ref{fig:final-loss} shows that the optimization process remains stable.
The validation curve follows the training curve closely, suggesting that the model does not simply memorize the training set.

\begin{figure}[H]
    \centering
    \maybeincludegraphics[width=0.62\linewidth]{../results/final/figures/true_vs_predicted.png}
    \caption{Analytical Black-Scholes prices against neural network predictions.}
    \label{fig:final-true-vs-predicted}
\end{figure}

Figure~\ref{fig:final-true-vs-predicted} confirms that predicted prices are tightly aligned with analytical prices across the observed price range.

\section{Monte Carlo Comparison}

The Monte Carlo benchmark is evaluated against the analytical Black-Scholes price on 512 deterministic test options.
The results are shown in Table~\ref{tab:mc-final-metrics}.

\begin{table}[H]
    \centering
    \begin{tabular}{lr}
        \toprule
        Metric             & Value    \\
        \midrule
        MAE                & 0.08882  \\
        RMSE               & 0.15790  \\
        $R^2$              & 0.999951 \\
        MAPE, prices $> 1$ & 0.6483\% \\
        \bottomrule
    \end{tabular}
    \caption{Final Monte Carlo metrics against analytical Black-Scholes prices.}
    \label{tab:mc-final-metrics}
\end{table}

The Monte Carlo benchmark also achieves high accuracy, as expected when the number of simulated paths is large.
However, it remains a simulation-based estimator and therefore has sampling error.
The neural network, in contrast, produces deterministic forward-pass predictions once trained.

\section{Feature Engineering Results}

An intermediate experiment compared the baseline feature set with the engineered moneyness feature.
Adding moneyness reduced MAE from 0.17269 to 0.15417 in the controlled intermediate setting with ReLU activations.
This suggests that exposing the ratio $S_0/K$ directly can help optimization, even though it does not add information beyond $S_0$ and $K$.

\section{Activation Function Results}

Another intermediate experiment compared ReLU, Tanh, and SiLU hidden activations using the same baseline feature set.
SiLU achieved the lowest MAE and RMSE among the tested activation functions.
This result is consistent with the smoothness of the Black-Scholes pricing function and motivated the final choice of SiLU.

\section{Runtime Results}

The runtime benchmark separates the initial neural network training cost from pricing time after training.
Using the improved final configuration, training required approximately 359 seconds in the benchmark environment.
Once trained, the neural network priced 1,000 options in approximately 0.0319 seconds on average.
The Monte Carlo benchmark with 20,000 paths required approximately 1.2824 seconds for the same number of options, while the vectorized analytical Black-Scholes formula required approximately 0.00091 seconds.

\begin{table}[H]
    \centering
    \begin{tabular}{lrr}
        \toprule
        Method                    & Mean time for 1,000 options & Mean options/s \\
        \midrule
        Black-Scholes formula     & 0.00091 s                   & 1,101,094.82   \\
        Neural network inference  & 0.03190 s                   & 31,350.77      \\
        Monte Carlo, 20,000 paths & 1.28235 s                   & 779.82         \\
        \bottomrule
    \end{tabular}
    \caption{Runtime benchmark for analytical, neural-network, and Monte Carlo pricing.}
    \label{tab:runtime-benchmark}
\end{table}

As expected, the analytical Black-Scholes formula is the fastest method when available.
The relevant surrogate-pricing comparison is therefore between neural network inference and Monte Carlo simulation.
In this benchmark, neural inference is roughly 40 times faster than Monte Carlo after the model has been trained.

\section{SVR Benchmark Results}

The reduced-scale SVR benchmark provides an additional comparison with a classical supervised regression model.
The experiment was repeated over three random seeds, using 5,000 synthetic contracts per run and the same engineered moneyness feature used by the final neural network.
Aggregate results are reported in Table~\ref{tab:svr-benchmark}.

\begin{table}[H]
    \centering
    \begin{tabular}{lrr}
        \toprule
        Metric             & Mean      & Standard deviation \\
        \midrule
        MAE                & 0.15088   & 0.00732            \\
        RMSE               & 0.23100   & 0.00592            \\
        $R^2$              & 0.999905  & 0.000005           \\
        MAPE, prices $> 1$ & 1.8253\%  & 0.2422\%           \\
        Fit time           & 12.3198 s & 0.7977 s           \\
        Prediction time    & 0.0301 s  & 0.0029 s           \\
        \bottomrule
    \end{tabular}
    \caption{Reduced-scale SVR benchmark aggregated over three random seeds.}
    \label{tab:svr-benchmark}
\end{table}

The SVR baseline also approximates the Black-Scholes pricing function well on reduced datasets.
This is indicated by its high $R^2$.
However, its absolute error is higher than that of the final neural network configuration.

\section{Noisy Target Robustness Results}

Table~\ref{tab:noisy-targets} reports the robustness experiment in which the neural network is trained on noisy Black-Scholes prices and evaluated against the clean analytical Black-Scholes target.

\begin{table}[H]
    \centering
    \begin{tabular}{lrrrr}
        \toprule
        Noise level & MAE     & RMSE    & $R^2$    & MAPE, prices $> 1$ \\
        \midrule
        0\%         & 0.10254 & 0.15169 & 0.999958 & 1.3213\%           \\
        1\%         & 0.11084 & 0.16105 & 0.999953 & 1.3960\%           \\
        5\%         & 0.17827 & 0.24167 & 0.999893 & 2.0943\%           \\
        \bottomrule
    \end{tabular}
    \caption{Noisy-target neural-network results evaluated against clean Black-Scholes prices.}
    \label{tab:noisy-targets}
\end{table}

The results show the expected degradation as label noise increases.
Nevertheless, the model remains accurate even under moderate target noise.

The same noisy-target experiment was also repeated with the reduced-scale SVR baseline.
Table~\ref{tab:noisy-svr-targets} reports aggregate metrics against clean Black-Scholes prices.

\begin{table}[H]
    \centering
    \begin{tabular}{lrrrr}
        \toprule
        Noise level & MAE     & RMSE    & $R^2$    & MAPE, prices $> 1$ \\
        \midrule
        0\%         & 0.15088 & 0.23100 & 0.999905 & 1.8253\%           \\
        1\%         & 0.19284 & 0.30765 & 0.999830 & 1.8705\%           \\
        5\%         & 0.56185 & 1.05057 & 0.998025 & 3.4903\%           \\
        \bottomrule
    \end{tabular}
    \caption{Noisy-target SVR benchmark evaluated against clean Black-Scholes prices.}
    \label{tab:noisy-svr-targets}
\end{table}

SVR remains accurate under mild label noise, but the degradation is stronger at the 5\% noise level.
Since the SVR benchmark is run on reduced datasets, this comparison is best interpreted as an additional classical robustness reference rather than as a direct replacement for the neural network experiment.

\section{Error Analysis}

\begin{figure}[H]
    \centering
    \maybeincludegraphics[width=0.72\linewidth]{../results/final/figures/error_distribution.png}
    \caption{Distribution of neural network prediction errors.}
    \label{fig:final-error-distribution}
\end{figure}

The error distribution in Figure~\ref{fig:final-error-distribution} is concentrated near zero, which is consistent with the low MAE and RMSE values reported in Table~\ref{tab:nn-final-metrics}.

\begin{figure}[H]
    \centering
    \begin{subfigure}{0.32\linewidth}
        \centering
        \maybeincludegraphics[width=\linewidth]{../results/final/figures/error_vs_moneyness.png}
        \caption{Moneyness}
    \end{subfigure}
    \begin{subfigure}{0.32\linewidth}
        \centering
        \maybeincludegraphics[width=\linewidth]{../results/final/figures/error_vs_maturity.png}
        \caption{Maturity}
    \end{subfigure}
    \begin{subfigure}{0.32\linewidth}
        \centering
        \maybeincludegraphics[width=\linewidth]{../results/final/figures/error_vs_volatility.png}
        \caption{Volatility}
    \end{subfigure}
    \caption{Absolute neural network error against selected financial features.}
    \label{fig:final-error-diagnostics}
\end{figure}

Figure~\ref{fig:final-error-diagnostics} is used to inspect whether the approximation error changes systematically across the parameter space.
This is important because aggregate metrics alone may hide localized weaknesses, especially around different moneyness regimes, short maturities, or high-volatility regions.
